\section{Introduction}
\label{sec:intro}

Concurrent systems rely on coordination mechanisms ranging from mutexes and
atomic instructions to database locks and distributed consensus protocols.
These mechanisms prevent concurrent actions from producing undesirable
outcomes, but add significant synchronization overheads; when coordination
crosses a network partition, it can also prevent availability.  Avoiding
unnecessary coordination is therefore a central goal in distributed systems and
databases.  This paper asks what coordination mechanisms accomplish at the
specification level, and when that effect is intrinsically required.

The CALM (Consistency As Logical Monotonicity) Conjecture
proposed a crisp, bidirectional answer: consistency can be achieved without
coordination exactly when the specification is monotone
~\cite{hellerstein2010declarative}.  Ameloot et al.\ proved it in both
directions, but only for Datalog queries evaluated by relational
transducer networks~\cite{ameloot2013relational,hellerstein2020keeping}. In a related vein, the CAP
Theorem proved that an available linearizable register cannot tolerate a
network partition~\cite{brewer2000towards,gilbert2002cap}.  These results are often
invoked more broadly than the formalisms support: CALM as a general
monotonicity boundary for coordination, CAP as a general constraint on distributed
systems.  The folklore expanded beyond what the literature proves.

Related results reach similar conclusions from different starting points.
Conflict-free Replicated Data Types (CRDTs)
use inflationary updates over join-semilattices to guarantee coordination-free
convergence~\cite{shapiro2011crdt}.  Highly Available Transactions (HATs) analyze
transaction histories to identify guarantees that remain available under
partition~\cite{bailis2013hat}.  Invariant confluence asks whether independently
reached, invariant-preserving database states can merge safely
~\cite{bailis2014coordination}.  Because these results speak different formal
languages, they do not translate directly.  Does CALM encompass CAP?  Are CRDT
join laws necessary for coordination freedom, or merely one strong sufficient
discipline?  Does convergent storage provide highly available transactions?
Do convergence and transaction availability preserve application invariants?

We develop a unified coordination criterion that spans this space,
including the translation questions above
(\S\ref{sec:cap-corollary}, \S\ref{sec:instances}).  We begin with a minimalist
specification model: a two-column relational table pairing input values with the
corresponding outputs that it permits.  The model is purely declarative: it
describes what can be observed about a system, without reference to any
imperative procedure.  The input domain $(I,\InputOrd)$ includes a partial order
describing how input grows; the output domain $(O,\OutputOrd)$ has its own
partial order characterizing when a later output preserves an earlier claim.
State-machine-style specifications can be compiled into this form
(\S\ref{sec:design-boundary}); the table elides operational details that are not
externally observable.  Despite its simplicity, the model captures both general
concurrency and consistency models that span applications---including
linearizability and causal consistency (\S\ref{app:registers})---and
application-specific requirements such as program invariants.  A
specification is \emph{consistent}---it has \emph{Consistency}---when the
executions it admits are \emph{confluent}: states (the current input together
with what can still happen) reached along different paths can still be brought
to a common state.  This unifies classical definitions of consistency: a specification
table captures the space of acceptable outcomes, while confluence asks whether
those requirements are maintained across concurrent executions.

We first focus on specification tables with a functional dependency from
input to output (\S\ref{sec:functional-specifications}). These \emph{functional
ordered specifications} capture arbitrary functions
\(f:(I,\InputOrd)\to(O,\OutputOrd)\). Although each input has one output, a system
may emit output before all input has arrived, and different executions may
receive input in different orders. We call this \emph{execution ambiguity}. A
function is monotone when \(i\InputOrd j\) implies \(f(i)\OutputOrd f(j)\): an
output emitted at \(i\) remains correct at \(j\).

Our first result is \emph{Functional CALM}
(\S\ref{sec:functional-calm-theorem}). It generalizes CALM to arbitrary
functions and partial orders, proving that Consistency is equivalent to
monotonicity.  When monotonicity fails, emitting some output requires
irrevocably ruling out later input that would contradict it, as an election
official does by closing the polls before announcing a result that later
ballots could overturn.  We call such an act a \emph{commitment}, and this kind
a \emph{seal}.  Functional CALM is a theorem about sealing: a functional
specification is consistent exactly when no output it permits needs a seal,
and a monotone specification can emit every output it permits without one.
This captures relational query languages and the usual CRDT state interface,
in which replica states are merged by a join.

We then give Consistency a distributed interpretation without changing the
specification model (\S\ref{sec:distributed-execution}). We place one local
execution at each participant of a distributed system, then set aside what does
not change specification state: message delivery, and each participant's
private logic for deciding what to do next (its \emph{strategy}).
\emph{Distributed Consistency} asks whether the resulting participant states can
still be brought to the same state. We prove that a specification has
Distributed Consistency if and only if it is consistent in the sense above.
For a consistent functional specification, state is then determined by input
alone, so participants with the same input agree regardless of delivery order
or emitted outputs: the convergence that strong eventual consistency asks of
replicated state, derived from the specification rather than from a merge
operator.

The functional specifications above cannot capture settings like leader election
or transaction ordering, where many outputs are permitted for a given input. In the
absence of a functional dependency determining the output for each input, a second source of
ambiguity arises, \emph{semantic ambiguity}, and with it a second job for
commitments: \emph{choosing}.  Picking one of two eligible leaders, or one
serial order for concurrent transactions, removes outputs that the input still
permits.  We show that sealing and choosing are the only two jobs: the effect
of any coordination mechanism on the specification is a combination of them.
To model semantic ambiguity, we must move beyond standard
monotonicity of functions. For the case of general relations we define an extension called
\emph{Lifted Monotonicity}, which asks whether outputs from different executions can remain
consistent with one another as input grows. Complete CALM proves that
Consistency, Lifted Monotonicity, and Distributed Consistency are equivalent
(\S\ref{sec:complete-calm-theorem}), so a specification is consistent exactly
when no output it permits needs a seal or a choice.  Both CALM theorems
characterize the same condition on the specification table: whenever permitted
rows have inputs with a common refinement $j$ (an input that extends all of
them), some output permitted at $j$
refines all of their outputs.  For one output per input this is ordinary
monotonicity; for many, it is Lifted Monotonicity.

Choice also explains mutual exclusion: when participants could make
conflicting choices, a system that requires Consistency must keep both from
completing, and coordination mechanisms are how it does so.  Consistency
itself guarantees only that a reconciling state always exists; reaching it is
up to the system, for instance by having every participant apply the same
tie-breaking rule (\S\ref{sec:complete-calm}).

Classical CALM tells us when coordination is unnecessary, but not whether a
chosen use of coordination is sufficient.  Complete CALM can answer the second
question because it represents commitments explicitly.  For example, total log
ordering and a key-value store are both non-monotone.  When they are composed
through a shared log order, however, committing to a single order is sufficient
to make the end-to-end specification monotone and consistent.  The key-value
store can then execute requests in the committed order without making further
commitments.  We call this analysis \emph{Proper Coordination} and formalize it
in Appendix~\ref{app:composition}.

Complete CALM also yields a Complete CAP Theorem
(\S\ref{sec:cap-corollary}).  Brewer conjectured that a partitioned system
cannot be both consistent and available; Gilbert and Lynch proved this for one
guarantee, the linearizable register.  Complete CAP makes the conjecture
precise for any consistency guarantee expressed as a specification, and in both
directions: it identifies which guarantees face the tradeoff and which escape
it.
A specification satisfying Lifted Monotonicity can keep its whole
\emph{interface}---every output it permits---available and confluent during a
partition; any other cannot.  A weaker \emph{response form} of Complete CAP
asks only that every operation be answered, where the interface
declares which outputs count as answers, and characterizes exactly when that is
possible (Appendix~\ref{app:cap-response}).  The Gilbert--Lynch theorem is the
special case in which one replica's only permitted answer is invalidated by
input it cannot see: answering would seal input that the other side of the
partition accepts.

This paper makes six contributions:
\begin{enumerate}
  \item \emph{A common specification model.} Ordered relational tables describe
        permitted input/output behavior; commitments record the irrevocable
        acts that remove alternatives from it, and each one seals input or
        chooses among outputs.
  \item \emph{Consistency as confluence.} One definition applies to general
        consistency models and application-specific requirements.
        This property of runs (sequences of input and output events) is shown
        to be equivalent to Distributed
        Consistency over distributed histories.
  \item \emph{Functional CALM.} For arbitrary functions between ordered
        domains, Consistency---and hence Distributed Consistency---is equivalent
        to monotonicity, and to no output needing a seal.
  \item \emph{Complete CALM and CAP.} With multiple outputs per input,
        Consistency---and hence Distributed Consistency---is equivalent to
        Lifted Monotonicity, the same common-refinement condition that
          underlies Functional CALM, and to no output needing a seal or a
          choice.  Complete CAP gives the corresponding
          partition boundary in two forms, for the whole interface and for
          answering every operation; the Gilbert--Lynch theorem is recovered as
          a special case.
  \item \emph{Composition and Proper Coordination.} Specifications compose by
        equijoin.  Committing a shared log order can make a non-monotone
        log/store composition consistent without further store commitments
        (Appendix~\ref{app:composition}).
  \item \emph{Coordination-free frontiers.} Two theorems characterize how much
        an interface can promise before coordination becomes necessary
        (Appendix~\ref{sec:frontiers}).
\end{enumerate}

Section~\ref{sec:functional-calm} introduces the framework on an election and
proves a chain of results, ending in the common-refinement characterization of
Consistency, that depend on output events only through a small \emph{event
contract}: four properties every output event satisfies.
Functional CALM then follows in a few lines.
Section~\ref{sec:distributed-execution} gives Consistency its distributed
reading.  Section~\ref{sec:ambiguous-specifications} lifts output events to
specifications with several outputs per input, verifies the same contract, and
shows that commitments only seal
or choose, so Section~\ref{sec:complete-calm}
assembles Complete CALM and Complete CAP from the same parts.
Section~\ref{sec:instances} applies the results to familiar systems, and
Section~\ref{sec:discussion} draws out the design lessons and scope limits.
Appendix~\ref{app:composition} proves the composition and Proper Coordination
results, Appendix~\ref{app:cap-response} develops the response form of
Complete CAP and recovers Gilbert--Lynch, Appendix~\ref{sec:frontiers} derives
maximal coordination-free
interfaces, and Appendix~\ref{app:additional-applications} gives further
classical applications.
